% 第 11 题图甲：探究弹簧振子周期与质量关系的装置（TikZ 重绘）
% 铁架台、弹簧、钩码（下表面吸附小磁体）、正下方的智能手机
\begin{tikzpicture}[font=\small, >=Stealth, scale=1.2]
  % 底座
  \draw[thick] (-0.2,0) -- (3.4,0) -- (3.3,0.22) -- (0,0.22) -- cycle;
  % 立柱
  \draw[thick] (0.45,0.22) -- (0.45,3.5);
  % 横臂
  \draw[thick] (0.3,3.25) -- (1.75,3.25);
  % 弹簧（折线）
  \draw[thick] (1.75,3.25) -- (1.75,3.15)
    -- (1.52,3.08) -- (1.98,2.96) -- (1.52,2.84) -- (1.98,2.72)
    -- (1.52,2.62) -- (1.75,2.55);
  % 钩码
  \draw[thick] (1.56,2.28) rectangle (1.94,2.55);
  % 小磁体
  \fill (1.58,2.14) rectangle (1.92,2.28);
  % 手机
  \draw[thick, rounded corners=1.5pt] (1.15,0.22) rectangle (2.35,0.52);
  \draw (1.3,0.29) -- (2.2,0.29);
  % 标注
  \node[left] at (0.38,1.7) {铁架台};
  \node[right] at (2.06,2.92) {弹簧};
  \draw[gray] (1.98,2.93) -- (1.9,2.9);
  \node[right] at (2.0,2.56) {钩码};
  \draw[gray] (1.99,2.52) -- (1.94,2.44);
  \node[right] at (2.0,1.98) {小磁体};
  \draw[gray] (1.99,2.0) -- (1.92,2.1);
  \node[right] at (2.4,0.4) {手机};
  \draw[gray] (2.36,0.4) -- (2.38,0.4);
  % 竖直虚线（磁体与手机的竖直对应关系）
  \draw[dashed, gray] (1.75,2.1) -- (1.75,0.52);
\end{tikzpicture}
